\chapter{Introduction - How to use this template}

\section{Choose your stile}

Since no style reference has been provided, we are quite free: you have first of all to choose the style and change some lines in the file \url{thesis.tex}. When two options are available, one of the two is commented starting the line with \%. To change it, comment what you don't need and de-comment the other option:

\subsection{Single- or Double-sided}
I like better the Single-sided, but if you want to go double-sided remove at line 14
\begin{verbatim}
\documentclass[12pt,a4paper,openright,oneside]{book}	%%One sided
\end{verbatim}
and de-comment line 15
\begin{verbatim}
\documentclass[12pt,a4paper,openright,twoside]{book}	%%Double sided
\end{verbatim}

\subsection{Headers and footers}
Lines 22-42 are going to use


\subsection{PDF settings}
Insert the information about you and the thesis in lines 56-61. They will be included in the PDF information
\begin{verbatim}
\hypersetup{
    pdfauthor={AuthorName},
    pdftitle={shortTitle},
    pdfsubject={subject},
    pdfkeywords={keyword1, keyword2}
}
\end{verbatim}

\subsection{Language}
The main language is the last, so for a thesis in english line 69 is
\begin{verbatim}
\usepackage[italian,english]{babel} 
\end{verbatim}
Swap ita with eng fi you are writing in italian as main language




\section[Sectioning]{Chapters, sections and subsections}

A thesis has three parts: head, body and tail. Three correspondent folders are containing the source files for these parts.

Any of this source files is imported in the \url{thesis.tex} file, which is the only one to get compiled. This is realised using the command
\begin{verbatim}
\input{subfolder/filename.tex}
\end{verbatim}
To exclude a section from the compiling process, comment or remove the correspondent line. If no subfolder is specified, \LaTeX\ looks for it in the same folder of \url{thesis.tex}

E.g. the introduction source file is in the body subfolder, and is included using the command \begin{verbatim} \chapter{Introduction - How to use this template}

\section{Choose your stile}

Since no style reference has been provided, we are quite free: you have first of all to choose the style and change some lines in the file \url{thesis.tex}. When two options are available, one of the two is commented starting the line with \%. To change it, comment what you don't need and de-comment the other option:

\subsection{Single- or Double-sided}
I like better the Single-sided, but if you want to go double-sided remove at line 14
\begin{verbatim}
\documentclass[12pt,a4paper,openright,oneside]{book}	%%One sided
\end{verbatim}
and de-comment line 15
\begin{verbatim}
\documentclass[12pt,a4paper,openright,twoside]{book}	%%Double sided
\end{verbatim}

\subsection{Headers and footers}
Lines 22-42 are going to use


\subsection{PDF settings}
Insert the information about you and the thesis in lines 56-61. They will be included in the PDF information
\begin{verbatim}
\hypersetup{
    pdfauthor={AuthorName},
    pdftitle={shortTitle},
    pdfsubject={subject},
    pdfkeywords={keyword1, keyword2}
}
\end{verbatim}

\subsection{Language}
The main language is the last, so for a thesis in english line 69 is
\begin{verbatim}
\usepackage[italian,english]{babel} 
\end{verbatim}
Swap ita with eng fi you are writing in italian as main language




\section[Sectioning]{Chapters, sections and subsections}

A thesis has three parts: head, body and tail. Three correspondent folders are containing the source files for these parts.

Any of this source files is imported in the \url{thesis.tex} file, which is the only one to get compiled. This is realised using the command
\begin{verbatim}
\input{subfolder/filename.tex}
\end{verbatim}
To exclude a section from the compiling process, comment or remove the correspondent line. If no subfolder is specified, \LaTeX\ looks for it in the same folder of \url{thesis.tex}

E.g. the introduction source file is in the body subfolder, and is included using the command \begin{verbatim} \chapter{Introduction - How to use this template}

\section{Choose your stile}

Since no style reference has been provided, we are quite free: you have first of all to choose the style and change some lines in the file \url{thesis.tex}. When two options are available, one of the two is commented starting the line with \%. To change it, comment what you don't need and de-comment the other option:

\subsection{Single- or Double-sided}
I like better the Single-sided, but if you want to go double-sided remove at line 14
\begin{verbatim}
\documentclass[12pt,a4paper,openright,oneside]{book}	%%One sided
\end{verbatim}
and de-comment line 15
\begin{verbatim}
\documentclass[12pt,a4paper,openright,twoside]{book}	%%Double sided
\end{verbatim}

\subsection{Headers and footers}
Lines 22-42 are going to use


\subsection{PDF settings}
Insert the information about you and the thesis in lines 56-61. They will be included in the PDF information
\begin{verbatim}
\hypersetup{
    pdfauthor={AuthorName},
    pdftitle={shortTitle},
    pdfsubject={subject},
    pdfkeywords={keyword1, keyword2}
}
\end{verbatim}

\subsection{Language}
The main language is the last, so for a thesis in english line 69 is
\begin{verbatim}
\usepackage[italian,english]{babel} 
\end{verbatim}
Swap ita with eng fi you are writing in italian as main language




\section[Sectioning]{Chapters, sections and subsections}

A thesis has three parts: head, body and tail. Three correspondent folders are containing the source files for these parts.

Any of this source files is imported in the \url{thesis.tex} file, which is the only one to get compiled. This is realised using the command
\begin{verbatim}
\input{subfolder/filename.tex}
\end{verbatim}
To exclude a section from the compiling process, comment or remove the correspondent line. If no subfolder is specified, \LaTeX\ looks for it in the same folder of \url{thesis.tex}

E.g. the introduction source file is in the body subfolder, and is included using the command \begin{verbatim} \chapter{Introduction - How to use this template}

\section{Choose your stile}

Since no style reference has been provided, we are quite free: you have first of all to choose the style and change some lines in the file \url{thesis.tex}. When two options are available, one of the two is commented starting the line with \%. To change it, comment what you don't need and de-comment the other option:

\subsection{Single- or Double-sided}
I like better the Single-sided, but if you want to go double-sided remove at line 14
\begin{verbatim}
\documentclass[12pt,a4paper,openright,oneside]{book}	%%One sided
\end{verbatim}
and de-comment line 15
\begin{verbatim}
\documentclass[12pt,a4paper,openright,twoside]{book}	%%Double sided
\end{verbatim}

\subsection{Headers and footers}
Lines 22-42 are going to use


\subsection{PDF settings}
Insert the information about you and the thesis in lines 56-61. They will be included in the PDF information
\begin{verbatim}
\hypersetup{
    pdfauthor={AuthorName},
    pdftitle={shortTitle},
    pdfsubject={subject},
    pdfkeywords={keyword1, keyword2}
}
\end{verbatim}

\subsection{Language}
The main language is the last, so for a thesis in english line 69 is
\begin{verbatim}
\usepackage[italian,english]{babel} 
\end{verbatim}
Swap ita with eng fi you are writing in italian as main language




\section[Sectioning]{Chapters, sections and subsections}

A thesis has three parts: head, body and tail. Three correspondent folders are containing the source files for these parts.

Any of this source files is imported in the \url{thesis.tex} file, which is the only one to get compiled. This is realised using the command
\begin{verbatim}
\input{subfolder/filename.tex}
\end{verbatim}
To exclude a section from the compiling process, comment or remove the correspondent line. If no subfolder is specified, \LaTeX\ looks for it in the same folder of \url{thesis.tex}

E.g. the introduction source file is in the body subfolder, and is included using the command \begin{verbatim} \input{body/introduction.tex} \end{verbatim}

\subsection{Front page}

Two front pages are provided:
\begin{enumerate}
	\item \url{frontPage.tex} for thesis with a single Relatore
	\item \url{frontPage-cr.tex} for thesis with Relatore and Correlatore
\end{enumerate}
choose what is needed adapting lines 94-95.

\subsection{Dedication}
Dedication is the initial inspirational quote. Can be edited in the file \url{dedication.tex} or removed.

\subsection{Chapter names}
Chapter and section names are in this form:
\begin{verbatim}
\chapter[Name in the index]{Name on the title}
\end{verbatim}
To exclude them from the numbering, use the *; If no name for the index is specified, will be assumed the same. For example:
\begin{verbatim}
\chapter*{Name on the title}
\end{verbatim}
won't be numbered and will have the same name in the title and in the index.

\section{Bulleting}

Bullet list example
\begin{itemize}
	\item first point
	\item second point
	\item third point \\
\end{itemize}

\noindent
Enumeration example
\begin{enumerate}
	\item first point
	\item second point
	\item third point\\
\end{enumerate}

\noindent
Description example
\begin{description}
	\item[first descr] first point
	\item[second descr]  second point
	\item[third descr]  third point\\
\end{description}

\noindent
\dots but you can also build nested lists
\begin{itemize}
	\item first point
	      \begin{itemize} %tab is only for the readability, but not necessary
		      \item first point
		      \item second point
	      \end{itemize}
	\item second point
	\item third point \\
\end{itemize}


\section{Floating objects}

Floating objects are tables, figures and so on.

\subsection{Position, captions and cross-referencing}

To force the position of an object use the options, for example:
\begin{verbatim}
\begin{figure}[htb]
\end{verbatim}
where [htb] means that the priority of positioning is \url{h} = here, \url{t} = top, \url{b} = bottom. Also \url{p} can be used and means that the object is placed in a page on itself.

Captions are managed automatically using the command
\begin{verbatim}
\caption{The title of my object}
\end{verbatim}
to cross reference to the object use the
\begin{verbatim} 
\label{labelName} 
\end{verbatim}
to assign a name to the object and
\begin{verbatim} 
\ref{labelName}
\end{verbatim}
to refer to that.



\subsection{Figures}

Insert a figure using the code
\begin{verbatim} 
\begin{figure}[htb]
\centering
\includegraphics[scale=0.15]{pictures/logo.png}
\caption{the logo of UniTo}
\label{myFigure}
\end{figure}
\end{verbatim}

then you can cross-reference to it. In this case the discussion is about Fig. \ref{myFigure}.

Where the number of the figure has been gotten using
\begin{verbatim} 
\ref{myFigure}
\end{verbatim}


\begin{figure}[htb]
	\centering

	\includegraphics[scale=0.15]{pictures/logo.png}
	\caption{the logo of UniTo}

\end{figure}

\subsection{Subfigures}

\begin{figure}[htb]
	\centering
	\subfigure[title of first subfigure]
	{\includegraphics[scale=0.15]{pictures/logo.png}}
	\subfigure[title of second subfigure]
	{\includegraphics[scale=0.3]{pictures/logo.png}}
	\caption{Two logos of UniTo of different sizes}

\end{figure}

\subsection{Tables}

 \end{verbatim}

\subsection{Front page}

Two front pages are provided:
\begin{enumerate}
	\item \url{frontPage.tex} for thesis with a single Relatore
	\item \url{frontPage-cr.tex} for thesis with Relatore and Correlatore
\end{enumerate}
choose what is needed adapting lines 94-95.

\subsection{Dedication}
Dedication is the initial inspirational quote. Can be edited in the file \url{dedication.tex} or removed.

\subsection{Chapter names}
Chapter and section names are in this form:
\begin{verbatim}
\chapter[Name in the index]{Name on the title}
\end{verbatim}
To exclude them from the numbering, use the *; If no name for the index is specified, will be assumed the same. For example:
\begin{verbatim}
\chapter*{Name on the title}
\end{verbatim}
won't be numbered and will have the same name in the title and in the index.

\section{Bulleting}

Bullet list example
\begin{itemize}
	\item first point
	\item second point
	\item third point \\
\end{itemize}

\noindent
Enumeration example
\begin{enumerate}
	\item first point
	\item second point
	\item third point\\
\end{enumerate}

\noindent
Description example
\begin{description}
	\item[first descr] first point
	\item[second descr]  second point
	\item[third descr]  third point\\
\end{description}

\noindent
\dots but you can also build nested lists
\begin{itemize}
	\item first point
	      \begin{itemize} %tab is only for the readability, but not necessary
		      \item first point
		      \item second point
	      \end{itemize}
	\item second point
	\item third point \\
\end{itemize}


\section{Floating objects}

Floating objects are tables, figures and so on.

\subsection{Position, captions and cross-referencing}

To force the position of an object use the options, for example:
\begin{verbatim}
\begin{figure}[htb]
\end{verbatim}
where [htb] means that the priority of positioning is \url{h} = here, \url{t} = top, \url{b} = bottom. Also \url{p} can be used and means that the object is placed in a page on itself.

Captions are managed automatically using the command
\begin{verbatim}
\caption{The title of my object}
\end{verbatim}
to cross reference to the object use the
\begin{verbatim} 
\label{labelName} 
\end{verbatim}
to assign a name to the object and
\begin{verbatim} 
\ref{labelName}
\end{verbatim}
to refer to that.



\subsection{Figures}

Insert a figure using the code
\begin{verbatim} 
\begin{figure}[htb]
\centering
\includegraphics[scale=0.15]{pictures/logo.png}
\caption{the logo of UniTo}
\label{myFigure}
\end{figure}
\end{verbatim}

then you can cross-reference to it. In this case the discussion is about Fig. \ref{myFigure}.

Where the number of the figure has been gotten using
\begin{verbatim} 
\ref{myFigure}
\end{verbatim}


\begin{figure}[htb]
	\centering

	\includegraphics[scale=0.15]{pictures/logo.png}
	\caption{the logo of UniTo}

\end{figure}

\subsection{Subfigures}

\begin{figure}[htb]
	\centering
	\subfigure[title of first subfigure]
	{\includegraphics[scale=0.15]{pictures/logo.png}}
	\subfigure[title of second subfigure]
	{\includegraphics[scale=0.3]{pictures/logo.png}}
	\caption{Two logos of UniTo of different sizes}

\end{figure}

\subsection{Tables}

 \end{verbatim}

\subsection{Front page}

Two front pages are provided:
\begin{enumerate}
	\item \url{frontPage.tex} for thesis with a single Relatore
	\item \url{frontPage-cr.tex} for thesis with Relatore and Correlatore
\end{enumerate}
choose what is needed adapting lines 94-95.

\subsection{Dedication}
Dedication is the initial inspirational quote. Can be edited in the file \url{dedication.tex} or removed.

\subsection{Chapter names}
Chapter and section names are in this form:
\begin{verbatim}
\chapter[Name in the index]{Name on the title}
\end{verbatim}
To exclude them from the numbering, use the *; If no name for the index is specified, will be assumed the same. For example:
\begin{verbatim}
\chapter*{Name on the title}
\end{verbatim}
won't be numbered and will have the same name in the title and in the index.

\section{Bulleting}

Bullet list example
\begin{itemize}
	\item first point
	\item second point
	\item third point \\
\end{itemize}

\noindent
Enumeration example
\begin{enumerate}
	\item first point
	\item second point
	\item third point\\
\end{enumerate}

\noindent
Description example
\begin{description}
	\item[first descr] first point
	\item[second descr]  second point
	\item[third descr]  third point\\
\end{description}

\noindent
\dots but you can also build nested lists
\begin{itemize}
	\item first point
	      \begin{itemize} %tab is only for the readability, but not necessary
		      \item first point
		      \item second point
	      \end{itemize}
	\item second point
	\item third point \\
\end{itemize}


\section{Floating objects}

Floating objects are tables, figures and so on.

\subsection{Position, captions and cross-referencing}

To force the position of an object use the options, for example:
\begin{verbatim}
\begin{figure}[htb]
\end{verbatim}
where [htb] means that the priority of positioning is \url{h} = here, \url{t} = top, \url{b} = bottom. Also \url{p} can be used and means that the object is placed in a page on itself.

Captions are managed automatically using the command
\begin{verbatim}
\caption{The title of my object}
\end{verbatim}
to cross reference to the object use the
\begin{verbatim} 
\label{labelName} 
\end{verbatim}
to assign a name to the object and
\begin{verbatim} 
\ref{labelName}
\end{verbatim}
to refer to that.



\subsection{Figures}

Insert a figure using the code
\begin{verbatim} 
\begin{figure}[htb]
\centering
\includegraphics[scale=0.15]{pictures/logo.png}
\caption{the logo of UniTo}
\label{myFigure}
\end{figure}
\end{verbatim}

then you can cross-reference to it. In this case the discussion is about Fig. \ref{myFigure}.

Where the number of the figure has been gotten using
\begin{verbatim} 
\ref{myFigure}
\end{verbatim}


\begin{figure}[htb]
	\centering

	\includegraphics[scale=0.15]{pictures/logo.png}
	\caption{the logo of UniTo}

\end{figure}

\subsection{Subfigures}

\begin{figure}[htb]
	\centering
	\subfigure[title of first subfigure]
	{\includegraphics[scale=0.15]{pictures/logo.png}}
	\subfigure[title of second subfigure]
	{\includegraphics[scale=0.3]{pictures/logo.png}}
	\caption{Two logos of UniTo of different sizes}

\end{figure}

\subsection{Tables}

 \end{verbatim}

\subsection{Front page}

Two front pages are provided:
\begin{enumerate}
	\item \url{frontPage.tex} for thesis with a single Relatore
	\item \url{frontPage-cr.tex} for thesis with Relatore and Correlatore
\end{enumerate}
choose what is needed adapting lines 94-95.

\subsection{Dedication}
Dedication is the initial inspirational quote. Can be edited in the file \url{dedication.tex} or removed.

\subsection{Chapter names}
Chapter and section names are in this form:
\begin{verbatim}
\chapter[Name in the index]{Name on the title}
\end{verbatim}
To exclude them from the numbering, use the *; If no name for the index is specified, will be assumed the same. For example:
\begin{verbatim}
\chapter*{Name on the title}
\end{verbatim}
won't be numbered and will have the same name in the title and in the index.

\section{Bulleting}

Bullet list example
\begin{itemize}
	\item first point
	\item second point
	\item third point \\
\end{itemize}

\noindent
Enumeration example
\begin{enumerate}
	\item first point
	\item second point
	\item third point\\
\end{enumerate}

\noindent
Description example
\begin{description}
	\item[first descr] first point
	\item[second descr]  second point
	\item[third descr]  third point\\
\end{description}

\noindent
\dots but you can also build nested lists
\begin{itemize}
	\item first point
	      \begin{itemize} %tab is only for the readability, but not necessary
		      \item first point
		      \item second point
	      \end{itemize}
	\item second point
	\item third point \\
\end{itemize}


\section{Floating objects}

Floating objects are tables, figures and so on.

\subsection{Position, captions and cross-referencing}

To force the position of an object use the options, for example:
\begin{verbatim}
\begin{figure}[htb]
\end{verbatim}
where [htb] means that the priority of positioning is \url{h} = here, \url{t} = top, \url{b} = bottom. Also \url{p} can be used and means that the object is placed in a page on itself.

Captions are managed automatically using the command
\begin{verbatim}
\caption{The title of my object}
\end{verbatim}
to cross reference to the object use the
\begin{verbatim} 
\label{labelName} 
\end{verbatim}
to assign a name to the object and
\begin{verbatim} 
\ref{labelName}
\end{verbatim}
to refer to that.



\subsection{Figures}

Insert a figure using the code
\begin{figure}[htb]
\centering
\includegraphics[scale=0.15]{pictures/logo}
\caption{the logo of UniTo}
\label{fig:myfigure}
\end{figure}



then you can cross-reference to it.
In this case the discussion is about Fig\ref{fig:myfigure}

Where the number of the figure has been gotten using
\begin{verbatim} 
\ref{myFigure}
\end{verbatim}


\begin{figure}[htb]
	\centering

	\includegraphics[scale=0.15]{pictures/logo.png}
	\caption{the logo of UniTo}

\end{figure}

\subsection{Subfigures}

\begin{figure}[htb]
	\centering
	\subfigure[title of first subfigure]
	{\includegraphics[scale=0.15]{pictures/logo.png}}
	\subfigure[title of second subfigure]
	{\includegraphics[scale=0.3]{pictures/logo.png}}
	\caption{Two logos of UniTo of different sizes}

\end{figure}

\subsection{Tables}

